\section{BinarySearch}
\paragraph{Ricorsione}Correttezza parziale di funzioni ricorsive via generazione delle chiamate di funzione sui cammini elementari:
\begin{itemize}
	\item Asserzioni in chiamata;
	\item Sommario di chiamata:
	\begin{lstlisting}[caption={profilo della funzione f}]
		§\aaa{@pre: F[p1,...,pn]}§ // sono i parametri formali di una funzione f
		§\aaa{@post: G[P1, ... Pn, rv]}§ // con rv valore restituito
		typeo f(type p1, ... type pn)   
	  //typeo $\color{Green}\omega$ := f(e1, ... ,en); $\color{Green}\rightarrow$ chiamata di f
	\end{lstlisting}
	\item Occorre almeno un cammino che termini con \texttt{@R: F[e1,...,en]} applicando la sostituzione di \texttt{p1,...,pn} con \texttt{e1,...,en}.
	\item Dimostrare la validità della condizione di verifica associata a ogni cammino che termina. L'asserzione di chiamata serve a dimostrare che la chiamata della funzione soddisfa la precondizione delle funzione.
	\item Occorre almeno un cammino che includa il sommario di chiamata \texttt{assume G[e1,...,en, v]; $\omega$ := v;} dove $v$ è il valore restituito dalla funzione chiamata ed \texttt{e1,...,en} sostituiscono \texttt{p1,...,pn}.
\end{itemize}
Nell'analisi di funzioni ricorsive, i cammini elementari iniziano con la precondizione e finiscono con la postcondizione o con un'altra asserzione.
\subsection{Esempio}
\begin{lstlisting}
	§\aaa{@pre: x >= 0}§
	§\aaa{@post: rv >= x}§
	int g(int x){...}
\end{lstlisting}
Considerando una funzione che chiama $g$ $$\texttt{$\omega$ := g(n+1)}$$
$n$ è definita nella funzione che chiama $g$. Bisogna annotare la chiamata di $g$:
\begin{lstlisting}
	§\aaa{@R1: n + 1 $\ge$ 0}§ // $\color{Green}\rightarrow$ asserzione in chiamata
	$\omega$ := g(n+1);
\end{lstlisting}
\begin{enumerate}
	\item Almeno un cammino che termina su $@R$
	\item Almeno un cammino che includa: \texttt{... assume v $\ge$ n+1; $\omega$ := v;} con \texttt{v} valore gestito da $g$.
\end{enumerate}